\documentclass[aspectratio=169,10pt]{beamer}

% 基本パッケージ
\usepackage{amsmath}
\usepackage{amssymb}
\usepackage{bm}
\usepackage{graphicx}

% カラーとテーブル用パッケージ
\usepackage{xcolor}
\usepackage{booktabs}
\usepackage{multirow}
\usepackage{multicol}
\usepackage{tcolorbox}

% 図とアルゴリズム用パッケージ
\usepackage{tikz}
\usepackage{algorithm}
\usepackage{algpseudocode}

% LuaTeX-jaの読み込み
\usepackage{luatexja}
\usepackage{luatexja-fontspec}

% フォントの基本設定
\usepackage{unicode-math}

% 数式フォントをセリフ体に設定
\setmathfont{Latin Modern Math}

\renewcommand{\kanjifamilydefault}{\gtdefault} % 漢字をゴシック体に指定
\usetheme[block=fill]{metropolis} % usepackageじゃなくてusethemeだから注意

% Proofブロックの定義
\newenvironment{proofblock}[1][]
{\begin{block}{#1}\vspace{0.5em}}
{\vspace{0.5em}\end{block}}

\newenvironment{outlineblock}[1][]
{\begin{tcolorbox}[colframe=black, colback=white, boxrule=1pt, boxsep=5pt]}
{\end{tcolorbox}}

% カラーの設定
\definecolor{titlecolor}{RGB}{0, 60, 120}
\definecolor{sectioncolor}{RGB}{0, 60, 120}
\definecolor{subsectioncolor}{RGB}{0, 60, 120}
\definecolor{boxcolor}{RGB}{240, 245, 255}
\definecolor{boxborder}{RGB}{200, 220, 240}
\definecolor{codecolor}{RGB}{240, 240, 240}

% beamerの設定
% \usetheme{Madrid}
% \usecolortheme{default}
\setbeamertemplate{navigation symbols}{}
\setbeamertemplate{footline}[frame number]  

\title{\textsc{全体講義}: 脅威モデルとセキュリティ}
\author{moz@sfc.wide.ad.jp}
\date{\today}

\begin{document}

\begin{frame}
\titlepage
\end{frame}

\begin{frame}{プロトコルとは？}
\begin{proofblock}[プロトコルの定義]
A set of rules governing the exchange or transmission of data between entities.
\end{proofblock}

% エンティティの定義を追加
\textbf{エンティティ}：プロトコルに参加する個々の主体（ユーザー、デバイス）を指す。



\begin{proofblock}[アルゴリズムの定義]
A set of instructions for performing a task or solving a problem.
\end{proofblock}

プロトコルは計算手順に焦点を当てるのに対し、プロトコルは参加者間の相互作用に重点を置く
一般的なプロトコルの例：HTTP、SMTP、TLS/SSL、SSH、OAuth、Paxos、Raft

\end{frame}

\begin{frame}{プロトコルの設計}
各プロトコルは特定の目的を達成するために設計される。
その目的を効果的に実現するためには、以下の点を明確に定義することが重要である：
\begin{itemize}
    \item どのようなエンティティがプロトコルに関与するか
    \item エンティティ間でどのような情報交換が行われるか
    \item \underline{それぞれのエンティティがどのような振る舞いをする可能性があるか}
\end{itemize}

\vspace{0.5em}

これらを明らかにして、\underline{安全なプロトコルを設計する鍵となる}。
\end{frame}


\begin{frame}{敵対者と脅威モデル}
プロトコル設計において、どのような振る舞いを許容するかを考える上では、\underline{エンティティの能力と行動に関する仮定を明確にする}ことが重要である。

\begin{proofblock}[敵対者]
An adversary refers to an attacker, typically with malicious intent, who undertakes actions against a system or protocol
敵対者は、プロトコルに参加するエンティティを制御下に置き、そのエンティティの動作を変更または監視できる状態。
\end{proofblock}

\vspace{0.5em}
\begin{proofblock}[脅威モデル]
\textcolor{red}{脅威モデル}（threat model）とは、敵対者の能力と行動に関する仮定を形式化したモデル。
\end{proofblock}

\end{frame}

\begin{frame}{脅威モデルの構成要素}
脅威モデルは以下の\underline{5つの要素から構成}から構成される：

\vspace{0.5em}

\begin{enumerate}
    \item \textbf{Corruption Model. }敵対者が制御するエンティティの行動様式
    \item \textbf{Computational Power. }敵対者が利用できる計算リソースの制約
    \item \textbf{Adaptivity. }敵対者がエンティティを腐敗させるタイミングと方法
    \item \textbf{Visibility. }敵対者が観察できる情報の範囲
    \item \textbf{Mobility. }敵対者によるエンティティの腐敗と回復のパターン
\end{enumerate}

\vspace{0.5em}

\begin{outlineblock}[]
この要素の組み合わせ方によって、\underline{異なる脅威モデル}が定義される。
どのような要素を選択するかは、プロトコルの目的、実行されうる環境によって決まる。
\end{outlineblock}
\end{frame}

\begin{frame}{腐敗モデル (Corruption Models)}
腐敗モデルは、敵対者はエンティティをどのように制御できるかを定義する。

\vspace{0.5em}

\begin{itemize}
    \item \textbf{Passive (Honest-but-Curious)}：プロトコルに完全に従いながら、観測したすべての情報を蓄積して活用できる。しかし、プロトコルの仕様に逸脱することはできない。\footnote{具体例として銀行の従業員、顧客の取引データを正しく処理しながらも、そのデータをこっそり記録して個人的に利用する場合。}
    
    \item \textbf{Crash}：応答を任意のタイミングで停止できるが、停止時間外は正確にプロトコルに従って動作する。
    
    \item \textbf{Omission}：エンティティは受信したのメッセージを任意に削除または無視できる。送信するメッセージの内容を変更することはできない。
    
    \item \textbf{Byzantine}：エンティティはプロトコル仕様から任意に逸脱できる。
\end{itemize}

\vspace{0.5em}

\begin{outlineblock}[]
各モデルは前のモデルを包含している：$\text{Byzantine} \supset \text{Omission} \supset \text{Crash} \supset \text{Passive}$
\end{outlineblock}

\end{frame}

\begin{frame}{適応性モデル (Adaptivity Models)}
適応性モデルは、敵対者がいつ、どのように参加者を腐敗させるかを定義する：

\vspace{0.5em}

\begin{itemize}
    \item \textbf{Static}：敵対者は腐敗させるエンティティをプロトコル実行前に選ぶ。実行中に変更できない。
    
    \item \textbf{\textcolor{red}{Adaptive}}：敵対者はプロトコル実行中に腐敗させるエンティティを動的に選択できる。
\end{itemize}

\vspace{0.5em}

\begin{outlineblock}[]
エンティティが少ない場合、プロトコルの通信フェーズ数が多い場合、プロトコルの実行時間が長い場合、Adaptiveモデルを考慮する。
例えば、鍵交換モデルなどでは、Staticモデルを考えることが一般的。
コンセンサスアルゴリズムでは、Adaptiveモデルを考えることが一般的。
\end{outlineblock}
\end{frame}


\begin{frame}{可視性モデル (Visibility Models)}
可視性モデルは、敵対者がどの程度のメッセージや状態（ノード内の情報）を観察できるかを定義する。

\vspace{0.5em}

\begin{itemize}
    \item \textbf{Full Information}：敵対者はすべてのメッセージ内容と各エンティティの内部状態（実行中の変数）を完全に観察可能
    
    \item \textbf{\textcolor{red}{Private Channels}}：正直なエンティティ間の通信内容は敵対者から隠されるが、通信の発生事実やメタデータ（送信元、宛先、タイミング、メッセージサイズ等）は観察可能。腐敗したエンティティの状態、メッセージ内容は観察可能。
    
    \item \textbf{Erasure Model}：Private Channelsに対して、参加者が腐敗する前に消去されたデータは、敵対者は観察・復元できない制約を足したもの。
    
    \item \textbf{Oblivious}：メッセージのヘッダー情報（送信元、宛先、長さ）のみを観察可能で、内容は完全に隠蔽される。\footnote{Trusted Computingのようなシステムでは、このモデルを考えることが多い。}
\end{itemize}

\begin{outlineblock}[]
    暗号プロトコルでは、機密性（intended to be kept secret）を保証するために、Private Channelsモデルを考えることが多い。
\end{outlineblock}

\end{frame}

\begin{frame}{計算能力モデル (Computational Power Models)}
計算能力モデルは、敵対者の計算資源の制限を定義する：


\begin{itemize}
    \item \textbf{Unbounded}：無制限の計算資源を持つ敵対者。理論上はすべての暗号を破ることができる。\\
    
    \item \textbf{Polynomially Bounded}：多項式時間の計算に制限される敵対者
    \footnote{
    ある問題(Problem)を解くときに、その問題の解を求める計算複雑性がオーダー$O(n^c)$ $n, c \in \mathbb{Z}$で表せる。\\
    例えば、離散対数問題、素因数分解問題の難しさなど。\\
    $Probabilistic$とは確率チューリングマシンにおいて、入力となるセキュリティパラメータ $\lambda \in \mathbb{Z}$（鍵長など）の大きさに対して， 乱数を用いながらも 多項式時間 $\lambda^c$以内で必ず終了するアルゴリズム}
    
    \item \textbf{Fine-grained}：特定の具体的な計算制限を持つ敵対者。\\
    秒間ハッシュ計算回数に制限があるモデル。Bitcoin等のProof-of-Workで使用
\end{itemize}

\vspace{0.5em}

\begin{outlineblock}[]
重要なのは、エンティティの計算能力ではなく、敵対者の計算能力である。
敵対者はVisibility Modelsにそって、エンティティのメッセージや内部状態を観察できる。その情報を敵対者はどのような計算資源を用いて活用できるかを定義する。
\end{outlineblock}
\end{frame}


\begin{frame}{移動性モデル (Mobility Models)}
移動性モデルは、敵対者による参加者の腐敗と回復のパターンを定義する：

\vspace{0.5em}

\begin{itemize}
    \item \textbf{Fixed}：一度腐敗した参加者はプロトコル実行中ずっと腐敗したままであり、回復することはない
    
    \item \textbf{\textcolor{red}{Mobile}}：敵対者は参加者を動的に腐敗させたり、回復させたりできる。ただし、任意の時点で腐敗している参加者の総数には上限がある
\end{itemize}

\vspace{0.5em}

\begin{outlineblock}[]
IoTやモバイルデバイス、システムの長期運用の視点を考えると、一般的にMobileモデルを考える。
Adaptivity ModelsのAdaptiveとの混同に注意。
\end{outlineblock}
\end{frame}

\begin{frame}{仮定の重要性}
セキュリティプロトコルの研究において、\textcolor{red}{仮定の明確化}は極めて重要である。

\vspace{0.5em}

\begin{itemize}
    \item 仮定が弱すぎる場合：現実的な攻撃に対して脆弱なプロトコルになる可能性がある
    \item 仮定が強すぎる場合：実装が複雑になり、実用性が低下する可能性がある
\end{itemize}

\vspace{0.5em}

仮定を明示する価値：
\begin{itemize}
    \item \textcolor{red}{セキュリティ保証の範囲}を正確に理解できる
    （「セキュリティ保証」とは、プロトコルが特定の脅威モデルの下で達成できる安全性の明確な指標）
    \item 異なるプロトコルを同じ基準で公正に比較できる
\end{itemize}

\vspace{0.5em}

実際のシステム設計では、特定の用途やリスク許容度に基づいて適切な仮定を選択することが重要である。
\end{frame}

\begin{frame}{例：Practical Byzantine Fault Tolerance (PBFT)}
\begin{proofblock}[PBFT]
\begin{itemize}
    \item PBFTは、\textcolor{red}{Byzantine障害を許容}しつつも、Linearlizableを達成するレプリケーションプロトコルである。
    \item ただし、$n \geq 3f+1$（ここで$f$は敵対者の数、$n$はレプリカの総数）のときに限る。
\end{itemize}
\end{proofblock}

PBFTプロトコルは、非同期通信モデルにおいて初めてLinearlizableを達成することを示したプロトコルである。
発表から20年経った現在でも、PBFTは依然として広く使われている。


\footnote{Linearlizableに複製するとは、すべてのレプリカが同じ順序でクライアント要求を処理することを意味する。各複製データにはシーケンス番号が割り当てられ、このシーケンス番号に従って処理される。}
\footnote{チューリング完全な状態マシンでは、複製されるデータをオペレーション（状態遷移関数）と見なす。オペレーション$f$, $g$がある場合、一般に$f \circ g \neq g \circ f \neq id$であるため、実行順序が重要になる。}
\end{frame}

\begin{frame}{PBFTのプロトコル詳細}
PBFTの主なフェーズ：
\vspace{0.5em}
\begin{enumerate}
    \item \textbf{Pre-prepare}：プライマリ（リーダー）$p$がクライアント要求$r$にシーケンス番号$s$を割り当て、$\langle \text{PRE-PREPARE}, s, r \rangle$メッセージを他のレプリカに送信
    
    \item \textbf{Prepare}：レプリカが要求の順序に合意し、$\langle \text{PREPARE}, s, r \rangle$メッセージを全レプリカに送信
    
    \item \textbf{Commit}：少なくとも$2f+1$のPREPAREメッセージを受信したレプリカは$\langle \text{COMMIT}, s, r \rangle$メッセージを送信し、実行を確約
    
    \item \textbf{Reply}：$2f+1$のCOMMITメッセージを受信したレプリカは要求を実行し、結果をクライアントに返送
\end{enumerate}
\end{frame}

\begin{frame}{PBFTの脅威モデル}
PBFTの脅威モデル：
\begin{proofblock}[PBFTの脅威モデル]
\begin{itemize}
    \item \textbf{Corruption}：\textcolor{red}{Byzantine} - レプリカは異なるメッセージを異なるレプリカに送信するなど、任意に振る舞うことができる
    \item \textbf{Adaptivity}：Adaptive - 敵対者はプロトコルの各フェーズで異なるレプリカを動的に腐敗させられる
    \item \textbf{Computational Power}：PPT - 多項式時間の計算に制限される（暗号署名が偽造不可能と仮定）
    \item \textbf{Visibility}：Private channels - 正直なレプリカ間の通信内容は秘密であり、内部状態も敵対者からは見えない
    \item \textbf{Mobility}：Mobile - 敵対者はプロトコルの各フェーズで異なるレプリカを動的に腐敗させられる
\end{itemize}
\end{proofblock}
\end{frame}

\begin{frame}{PBFTのセキュリティ保証}
PBFTは以下の重要なセキュリティ特性を提供する：

\begin{proofblock}[PBFTの脅威モデル]
$n \geq 3f + 1$のとき（$f$は敵対者の数、$n$はレプリカの総数）：
\begin{itemize}
    \item \textbf{\textcolor{red}{安全性} (Safety)}：正直なレプリカが同じシーケンス番号で異なる値をコミットすることはない
    \item \textbf{\textcolor{red}{活性} (Liveness)}：すべてのクライアント要求は最終的に処理され応答が返される
\end{itemize}
\end{proofblock}

\vspace{0.5em}

脅威モデルの条件において、SafetyとLivenessを両方満たす場合、PBFTはセキュアであると言える。

\vspace{0.5em}


\textbf{Question}: なぜ$n \geq 3f + 1$のときSafetyが満たされるのか？ PBFTがセキュアであることを説明してみよう！
\end{frame}

\begin{frame}{PBFTのSafety証明}
仮定: 二つの異なる値$v$と$v'$が同じシーケンス番号$s$でコミットされたと仮定する。

\begin{lemma}
各値がコミットされるためには、少なくとも$2f+1$個のレプリカからPREPAREメッセージが必要である。
\end{lemma}

\begin{itemize}
    \item 値$v$がコミットされるには少なくとも$2f+1$個のレプリカからのPREPAREメッセージが必要
    \item 同様に、値$v'$がコミットされるには少なくとも$2f+1$個のレプリカからのPREPAREメッセージが必要
    \item 合計$n \leq 3f+1$レプリカのうち、悪意あるレプリカは最大$f$個
    \item $2f+1 + 2f+1 - n \geq 2f+1 + 2f+1 - (3f+1) = f+1$
    \item つまり、少なくとも$f+1$個の\textcolor{red}{正直なレプリカ}が両方の値$v$と$v'$についてPREPAREメッセージを送信したことになる
    \item しかし、正直なレプリカは同じシーケンス番号に対して複数の値にPREPAREを送信しない（プロトコルの定義により）
    \item これは矛盾であり、したがって仮定は誤り
\end{itemize}
\end{frame}

\begin{frame}{合成可能性問題 (Composability Problem)}
ちょっとお題が変わって、複雑なプロトコルの証明について考えてみよう。

分散システムでは、既存のプロトコル上に新しいプロトコルを構築することが一般的である\footnote{例えば、PBFTを使用するHyperledger Fabric上に新たなアプリケーションを構築}

\vspace{0.5em}

\textbf{Motivation}：複数のプロトコル$\pi_1, \pi_2, \cdots, \pi_n$があったとき、それらを組み合わせたプロトコル$\Pi^{\pi_1, \pi_2, \cdots, \pi_n}$の証明はどう行うか？

システム全体を最初から再証明するのは非常に複雑で時間がかかる。
\end{frame}

\begin{frame}{Universal Composability Framework}
解決策の一つ：\textcolor{red}{Universal Composability (UC)}を用いたプロトコルの安全性証明
セキュリティ分析において考慮される攻撃の範囲と限界を明示的に定義する。明示的な脅威モデリングが不可欠である。
\vspace{0.5em}

UCは、プロトコルを部品の集合として捉え、これらの部品を合成したプロトコルのセキュリティを保証する。
これにより、複雑なシステムの分析を、より単純な構成要素の分析に還元できる。
\end{frame}

\begin{frame}{Universal Composability Framework}
Universal Composability (UC)の中心的なアイデアは、\textcolor{red}{理想/現実世界パラダイム}である：

\vspace{0.5em}

\textbf{理想世界 (Ideal World)}：
    \begin{itemize}
        \item 理想機能 $\mathcal{F}$ が直接タスクを安全に実行する
        \item これは「定義上安全」な仕様を表す抽象的なエンティティ
        \item 例：理想的な署名機能 $\mathcal{F}_{\text{sig}}$ は、正当な署名者のみが署名を生成でき、偽造が不可能であることを保証
    \end{itemize}
     
\vspace{0.5em}

\textbf{現実世界 (Real World)}：
    \begin{itemize}
        \item 実際のプロトコル $\Pi$ が参加者によって実装され、敵対者 $\mathcal{A}$ の攻撃に晒される
        \item 敵対者は、前述の脅威モデル（腐敗、適応性、計算能力等）に基づいて定義される
        \item 攻撃に晒されるということは、理想世界のような「定義上の安全性」が自明でなく、証明が必要になる
    \end{itemize}
\end{frame}

\begin{frame}{Universal Composability Framework（続き）}
\textbf{環境} $\mathcal{Z}$：
    \begin{itemize}
        \item プロトコル外部のすべてのシステム（ユーザー、他のプロトコル等）を表現する
        \item 入力を提供し（例：署名すべきメッセージ）、出力を受け取る（例：署名された結果）
        \item 環境はプロトコル参加者と敵対者の両方と相互作用し、どちらの世界（理想/現実）で動作しているかを区別しようとする
    \end{itemize}
    
\textbf{\textcolor{red}{シミュレータ}} $\mathcal{S}$：
    \begin{itemize}
        \item 理想世界で敵対者を模倣するエンティティ
        \item 現実世界の敵対者の振る舞いを理想世界でシミュレートし、環境から見て区別不可能にする
    \end{itemize}

\begin{proofblock}[セキュリティの定義]
任意の敵対者 $\mathcal{A}$ に対して、シミュレータ $\mathcal{S}$ が存在し、どんな環境 $\mathcal{Z}$ も以下を区別できないこと：
\end{proofblock}
\begin{itemize}
    \item 実際のプロトコル $\Pi$ と敵対者 $\mathcal{A}$ との対話
    \item 理想機能 $\mathcal{F}$ とシミュレータ $\mathcal{S}$ との対話
\end{itemize}

\end{frame}

\begin{frame}{本当に仮定はこれだけでいいのか？}
ここまでで用いた仮定は以下の通りである：
\begin{itemize}
    \item \textbf{ネットワークモデル}：非同期か同期か、部分同期（通信遅延に上限がある）か
    \item \textbf{敵対者モデル}： 敵対者はどのように振る舞うか
    \item \textbf{腐敗ノードの数}：全体の何割まで敵対者が制御できるか ($n \geq 3f+1$ など)
\end{itemize}

\textbf{Question}: 実世界で動かす上で考えるべき脅威、仮定はこれだけか？
\end{frame}

\begin{frame}{Sybil Attack}
これだけではすべてのプロトコルにSybil Attackが適用できてしまう。

\begin{proofblock}[Sybil Attack]
単一の実体が複数の偽の身元を作成し、システム内で不当な影響力を得る攻撃である。
\end{proofblock}

例：PBFTでは$n \geq 3f+1$の条件がある。$n$はノードの数、$f$は敵対者の数である。
敵対者が単一のエンティティ（敵対者）が複数のエンティティとして振る舞えることを許すと、容易に条件を破ることができる。
\end{frame}

\begin{frame}{アイデンティティの生成コスト}
公開鍵を識別子として用いる場合、大抵の計算機では公開鍵は計算コストが低いため、Sybil Attackに脆弱である。

単にノードが公開鍵に対応する秘密鍵を持っていることをデジタル署名を用いて示すだけでは、そのノードが独立した実体\footnote{エンティティをどのように区別し、何を持って独立と定義するのかは議論の余地がある。}であることを保証できない。

単一のエンティティと公開鍵の一対一対応を保証するためには、認証局による一対一対応の検証と証明が必要である。
\begin{itemize}
    \item 認証局の存在を許す = 信頼できる第三者の存在を仮定する
    \item 単一障害点とスケーラビリティの問題が生じる\footnote{あるエンティティが自分は山田花子と主張したが、本当に山田花子であるかどうかを検証したあとに、山田花子とそのエンティティが提示した公開鍵を紐付ける。その紐付きに対して認証局か署名を行い、証明書を作成する。}
\end{itemize}

ほかにも\textbf{Proof-of-Work/Stake}：新たな識別子の作成に計算資源や経済的資源を要求する研究もある。
\end{frame}

\end{document}